\section{Conclusion}

EthnoVerse presents a comprehensive approach to digital cultural preservation by integrating a scalable archival platform with immersive 3D visualization technologies. The system successfully combines full-stack web development, structured metadata management, and modern 3D reconstruction techniques to create an interactive and accessible repository for marginalized cultural knowledge.

A key contribution of this project is the integration of 3D scene reconstruction into a cultural archive. While Neural Radiance Fields (NeRF) were initially explored for their high-fidelity rendering capabilities, the transition to 3D Gaussian Splatting (3DGS) proved to be a critical design decision. 3DGS enabled significantly faster training and real-time rendering, making it suitable for web-based deployment and interactive user experiences. This shift highlights the importance of aligning technical choices with practical constraints such as performance, scalability, and accessibility.

The system demonstrates that cultural archives need not remain static repositories; instead, they can evolve into dynamic, experiential platforms that allow users to engage with cultural environments in meaningful ways. By linking media, transcripts, and reconstructed spaces, EthnoVerse provides a richer, multi-modal representation of community knowledge.

Beyond its technical contributions, the project emphasizes ethical considerations in digital archiving, including informed consent, contextual metadata, and responsible representation. These aspects are essential when working with historically underrepresented communities.

Overall, EthnoVerse establishes a strong foundation for future work at the intersection of computer science, digital humanities, and cultural preservation.

\section{Future Work}

While the current system achieves its core objectives, several avenues exist for further enhancement and expansion.

\subsection{Scalability and Multi-Community Support}

Future iterations of the system can extend beyond a single community to support multiple cultural groups. This would require improved data organization, scalable storage solutions, and more robust indexing mechanisms to handle larger datasets.

\subsection{Mobile-Based Data Collection}

Developing a mobile application for data collection would enable field researchers and community members to capture and upload media directly. This would streamline the data ingestion process and encourage participatory archiving.

\subsection{Enhanced 3D Reconstruction}

Although 3D Gaussian Splatting enables real-time rendering, further improvements can be explored:
\begin{itemize}
    \item Hybrid approaches combining 3DGS with mesh extraction for better geometry
    \item Improved handling of dynamic scenes and moving subjects
    \item Optimization for low-resource devices such as mobile browsers
\end{itemize}

\subsection{User Interaction and Immersive Features}

Future work can enhance user engagement by introducing:
\begin{itemize}
    \item Guided virtual tours within reconstructed environments
    \item Annotation layers for cultural storytelling
    \item Virtual or augmented reality (VR/AR) support
\end{itemize}

\subsection{AI-Assisted Metadata and Search}

The system can be extended using advanced AI techniques:
\begin{itemize}
    \item Automatic tagging of images and videos using computer vision
    \item Semantic search using embeddings instead of keyword matching
    \item Multilingual transcription and translation support
\end{itemize}

\subsection{Community Participation and Ethics}

A key direction for future work is enabling community-driven contributions:
\begin{itemize}
    \item User-submitted content with moderation workflows
    \item Consent management dashboards
    \item Tools for community members to curate and narrate their own archives
\end{itemize}

These improvements would not only enhance the technical robustness of EthnoVerse but also strengthen its role as a participatory and ethically grounded cultural preservation platform.